\section{An explicit uniform bound from the eightfold supplier construction}
\label{sec:eightfold-quantitative}

The following quantitative extraction is new to this run. It uses the geometric supplier and dominance inputs of \citet{Severn26}, with the manuscript status and independent checks recorded in Section~\ref{sec:eightfold}. It does not attribute the numerical bound below to that manuscript. The construction on powers of quotient curves comes from \citet[Theorem~2.0 and Lemmas~2.3--2.6]{Schoen88}, with Severn's local repair where required. The invariant-theory input is the linear-algebraic part of \citet[Lemma~3.5, Proposition~3.6, and Remark~3.7]{Milne99}. Every subsequent denominator and degree estimate is proved below.

We keep $K=\Q(\omega)$ and $\cO=\Z[\omega]$. In this section a polarization class is denoted by $\theta$. For an integral signed cycle $Z=\sum_j a_jZ_j$ of pure dimension $d$, define its degree mass by
\[
 \operatorname{mass}_\theta(Z)=\sum_j|a_j|\int_{Z_j}\theta^d.
\]
It bounds the degrees of its effective positive and negative parts separately. All degrees remain unnormalized intersection degrees.

\begin{theorem}[SQ1: an explicit effective seed degree]\label{SQ1}
Consider a fixed polarized split Weil eightfold family for $K$ whose integral homology lattice is $\cO$-stable. Let $\Delta$ be the product of the elementary divisors of its polarization. Define the following positive integers:
\begin{align}
 G&=39\,191\,040, & g&=4G+1=156\,764\,161,\label{eq:SQ-Gg}\\
 R&=(2G+1)G!G^G, & L&=2gR,\label{eq:SQ-RL}\\
 S_0&=G^{2g}(g!)^2, & S_B&=(L/2)!S_0^R,\label{eq:SQ-S}\\
 k&=L/2+4, & B_0&=(40G)^k\binom{k}{4}S_B\,8!\,G^8,\label{eq:SQ-B}\\
 N_{\rm isog}&=4G^8\Delta, & D_8(\Delta)&=N_{\rm isog}^{8}B_0.\label{eq:SQ-D}
\end{align}
Here $L$ in \eqref{eq:SQ-RL} is an integer bound for the number of curve factors, not a line bundle.

On a nonempty analytic open subset of this family, every fibre admits an effective integral four-cycle $Z$ with
\[
 \deg_\theta Z\le D_8(\Delta),\qquad P_4\cl(Z)\ne0,
\]
where $P_4$ is the Weil projector of \findingref{P1}. More precisely, there is a finite collection of effective four-cycles of this degree bound whose projected classes span the Weil plane. Consequently \findingref{P3} and \findingref{R21} give explicit uniform bounded division on the whole fixed family. Formulas for the final bound are given in \findingref{SQ9} below.
\end{theorem}

The factorial expression is an explicit integer, despite its enormous size. It is neither a practical complexity estimate nor a claim about the least possible multiplier or cycle degree. The proof of the theorem is the sequence \findingref{SQ2}--\findingref{SQ7}.

\subsection{Integral projectors and polarization degree}

\begin{proposition}[SQ2: a common projector denominator]\label{SQ2}
Let $\mathcal G$ be the finite reflection group reconstructed from Severn's printed roots, so $|\mathcal G|=G$. A cover $C\to\mathbf P^1$ with the printed twenty-involution tuple has genus $g=4G+1$. Its eightfold factor $A\subset J(C)$ can be defined by a Rosati-symmetric primitive idempotent $e$ with
\[
 Ge\in\Z[\mathcal G].
\]
The lattice $H_1(A,\Z)$ is $\cO$-stable. If $\theta_A$ is the restriction of the canonical principal polarization of $J(C)$, each elementary divisor of $\theta_A$ divides $G$. In particular,
\begin{equation}\label{eq:SQ-sourcepolarization}
 \Delta_A\mid G^8,\qquad \int_A\theta_A^8=8!\Delta_A\le8!G^8.
\end{equation}
The integer $G$ also clears the primitive rational projectors of all cyclic suppliers.
\end{proposition}
\begin{proof}
There are twenty branch points with inertia order two. Thus Riemann--Hurwitz gives
\[
 2g(C)-2=G\left(-2+20\left(1-\frac12\right)\right)=8G.
\]
This uses twenty involutions, not ten elements of order three.

Let $\sigma$ be the six-dimensional representation over $K$, let $\chi$ be its character, and choose a printed reflection $s$. The rational central idempotent for the conjugate pair $\sigma,\bar\sigma$ is
\[
 \varepsilon=\frac6G\sum_{a\in\mathcal G}
       \operatorname{Tr}_{K/\Q}\bigl(\chi(a^{-1})\bigr)a.
\]
The trace in every coefficient is a rational algebraic integer and hence belongs to $\Z$. The representation is realizable over $K$, so its rational simple factor is $M_6(K)$. Set
\[
 e=\varepsilon\frac{1-s}{2}.
\]
On that factor this is the orthogonal projection onto the one-dimensional negative eigenspace of the reflection. It is therefore a primitive idempotent. Both $\varepsilon$ and $s$ are fixed by the Rosati involution $a\mapsto a^{-1}$, and they commute, so $e$ is Rosati symmetric. The factor $6/2=3$ shows directly that $Ge\in\Z[\mathcal G]$. Its image in the Jacobian has dimension eight because the multiplicity space of $\sigma$ has $K$-dimension eight, as in \citet[Section~3, Definition \texttt{def:pt} and Lemma \texttt{lem:weiltype}]{Severn26}.

The exact group certificate also verifies that all six scalar units of $K$ belong to $\mathcal G$. In particular the scalar $\omega I_6$ is a central group element. It commutes with $e$, preserves the abelian subvariety $A$, and acts there with minimal polynomial $X^2+X+1$. It consequently preserves $H_1(A,\Z)$ and supplies the asserted integral $\cO$-action. This source-side integrality will be used in \findingref{SQ7}.

Write $\Lambda=H_1(J(C),\Z)$ and let $E$ be its unimodular alternating form. The lattice of $A$ is the primitive sublattice
\[
 \Lambda_A=\Lambda\cap e\Lambda_\Q.
\]
Its dual lattice for $E|_{e\Lambda_\Q}$ is $\Lambda_A^\#=e\Lambda$. To prove this, first note that $E(ex,y)=E(x,y)$ for $x\in\Lambda$ and $y\in\Lambda_A$, by self-adjointness of $e$. Conversely, every integral functional on the primitive sublattice $\Lambda_A$ extends to $\Lambda$. Unimodularity represents that extension by pairing with an element of $\Lambda$, and projection by $e$ gives the required representative in $e\Lambda$. Hence
\[
 G\Lambda_A^\#=Ge\Lambda\subseteq\Lambda_A.
\]
The polarization kernel $\Lambda_A^\#/\Lambda_A$ is annihilated by $G$. Its elementary divisors therefore divide $G$, proving \eqref{eq:SQ-sourcepolarization}.

Finally, for a supplier $(H',\chi')$ with $m=\operatorname{ord}(\chi')$, its primitive rational projector has coefficients
\[
 \frac1{|H'|}\sum_{t\in(\Z/m)^\times}\chi'(a)^{-t},\qquad a\in H'.
\]
The sums are integral Ramanujan sums, and $|H'|\mid G$. Thus $G$ clears all these projectors as well.
\end{proof}

\subsection{An explicit bound for the supplier identity}

\begin{proposition}[SQ3: bounded integral supplier coefficients]\label{SQ3}
For the printed reflection branch data there exists an identity of determinant characters
\[
 e_\sigma=\sum_i n_i a(\chi_i),\qquad
 \sum_i|n_i|\le R=(2G+1)G!G^G,
\]
using admissible suppliers $(H'_i,\chi_i)$ with $m_i=\operatorname{ord}(\chi_i)\ge3$ and positive even primitive multiplicity $h_i$. This bound does not depend on the optional deposited $84$-row identity.
\end{proposition}
\begin{proof}
Brauer induction expresses $\sigma$ as an integral combination of representations induced from linear characters of subgroups. Applying the determinant-character homomorphism gives the required membership in the lattice of supplier characters; this is the proof of \citet[Section~5, Theorem \texttt{thm:brauer}]{Severn26}. We recall why its discarded terms contribute zero.

A character of order one or two is real, so the multiplicities for a nonreal pair $\tau,\bar\tau$ agree. Its free determinant coordinates vanish. For a quaternionic $\tau$, the representation carries a nondegenerate invariant alternating form. The isotypic summand of a real one-dimensional character is orthogonal to every other isotypic summand, and the restricted form is nondegenerate. Its multiplicity is therefore even, so the coordinate modulo two also vanishes. If $m\ge3$ but $h=0$, the identity
\[
 h=\sum_\tau m_\tau
       \operatorname{mult}_\chi(\tau|_{H'})
\]
is a sum of nonnegative integers. Every multiplicity with $m_\tau>0$ is zero, and the determinant column is again zero. For every remaining term, a nontrivial cyclic branch exponent is $m/2$, because the original inertia is an involution. The cyclic surface-group relation makes the number of these exponents even. They thus pair with their negatives, and $h=2q-2+s$ is even. This proves admissibility of all retained terms, including the unramified case.

The character group has the form $\Z^a\oplus(\Z/2)^b$. Append columns $2e_j$ for the torsion coordinates, and represent the problem in $\Z^{a+b}$. The target has coordinates of absolute value at most one. Every supplier column has entries of absolute value at most $G$: restriction multiplicities are at most the dimensions of the corresponding irreducible representations, and these dimensions are at most $G$. There are at most $G$ irreducible representations, so the number of rows is at most $G$.

We give the integer-linear-algebra bound explicitly. Let an integer column matrix have rank $r$, entries bounded by $M\ge1$, and let $b$ be a vector of coordinate size at most one in its integral column lattice. Choose $r$ independent rows and $r$ independent columns, forming an invertible square matrix $B$. Projection to the selected rows is injective on the rational column span. Set
\[
 d=|\det B|\le r!M^r.
\]
The subgroup of $\Z^r/B\Z^r$ generated by the other columns has at most $d$ elements. In its Cayley graph a shortest path from zero to the class of $b$, using those columns and their negatives, has at most $d-1$ steps. After subtracting those signed columns, the residual vector $v$ lies in $B\Z^r$ and satisfies
\[
 \|v\|_\infty\le1+(d-1)M\le dM.
\]
Cofactors of $B$ have absolute value at most $(r-1)!M^{r-1}$. Cramer's rule, including its denominator $d$, gives
\[
 |(B^{-1}v)_j|\le r!M^r.
\]
The total absolute coefficient is therefore at most $(r+1)r!M^r$. With $r\le G$ and $M=G$, this is bounded by $(2G+1)G!G^G$. Removing the auxiliary columns for the order-two coordinates cannot increase the supplier coefficient sum. This proves the displayed bound.
\end{proof}

\subsection{Degrees of Schoen's actual cycles}

\begin{proposition}[SQ4: a bounded integral cycle on each curve power]\label{SQ4}
For an admissible supplier let $C'=C/\ker\chi$, let $m=\operatorname{ord}(\chi)$, and let $h>0$ be its primitive multiplicity. On $C'^h$ there is an integral signed cycle $Y$ of dimension $h/2$, whose class belongs to the determinant subspace of $H^h(C'^h,\Q)$ and has nonzero component in every primitive determinant eigenline. If $H$ is the sum of the degree-one point classes on its factors, then
\begin{equation}\label{eq:SQ-singlemass}
 \operatorname{mass}_H(Y)\le m^{3h/2}(h!)^2
                  \le G^{2g}(g!)^2=S_0.
\end{equation}
\end{proposition}
\begin{proof}
The rational primitive summand of $H^1(C',\Q)$ has dimension $\varphi(m)h$. Since $m\ge3$, $\varphi(m)\ge2$, and therefore
\[
 h\le g(C')\le g(C)=g.
\]
Also $m\le G$. These observations justify the final estimate in \eqref{eq:SQ-singlemass} once the first estimate is proved.

Put $X=C/H'$. Schoen's construction uses the finite quotient
\[
 q:C'^h\longrightarrow W_0=C'^h/G_0,
 \qquad \beta_0:W_0\longrightarrow\operatorname{Sym}^hX,
 \qquad |G_0|=m^{h-1}h!,
\]
where $G_0$ is the kernel of the sum map on $(\Z/m)^h$, extended by the symmetric group. It constructs a linear space $P_0\simeq\mathbf P^{h/2}$ in an Abel--Jacobi fibre. The prescribed blow-ups leave its strict transform isomorphic to $P_0$, and the cyclic cover splits over it. A component pushes to $Q_0\subset W_0$ of degree one over $P_0$. The normal-bundle calculation and exceptional-divisor argument give nonzero primitive character components of $[Q_0]$ \citep[Theorem~2.0, Lemmas~2.3--2.6, pp.~11--17]{Schoen88}; for the exceptional exponents one uses \citet[Section~4, Lemma \texttt{lem:schoenR}]{Severn26}.

For a general point $x\in X$, the incidence divisor
\[
 D_x=x+\operatorname{Sym}^{h-1}X
\]
restricts to the hyperplane class on $P_0$. Indeed, on the complete linear system constituting the Abel--Jacobi fibre, containing $x$ is one linear condition; restriction to the chosen linear subspace gives $\cO_{P_0}(1)$ for general $x$. Consequently
\[
 \int_{Q_0}(\beta_0^*D_x)^{h/2}=1.
\]
The pullback of $D_x$ to $C'^h$ is numerically $mH$.

Choose an irreducible component $Z_0$ of $q^{-1}(Q_0)$ and let $e_0$ be its degree over $Q_0$. Since $q$ is finite, $e_0\le|G_0|$. The projection formula gives
\[
 \deg_H Z_0=\frac{e_0}{m^{h/2}}.
\]
To keep the construction integral without assuming flatness of $q$ along $Q_0$, use the orbit sum
\[
 \mathcal R_0=\sum_{a\in G_0}a_*Z_0,
\]
retaining repetitions. The finite-quotient transfer identity
\[
 q^*q_*\cl(Z_0)=\sum_{a\in G_0}a_*\cl(Z_0)
\]
shows that $\cl(\mathcal R_0)=e_0q^*[Q_0]$. This is an identity in rational cohomology; it follows equally from the transfer on the finite quotient and Poincar\'e duality for a quotient of a smooth variety. The cycle $\mathcal R_0$ itself is an integral effective cycle. Its degree is bounded by
\[
 \deg_H\mathcal R_0
  =\frac{|G_0|e_0}{m^{h/2}}
  \le\frac{|G_0|^2}{m^{h/2}}.
\]

Let $\delta(t)$ be the lift of the residual cyclic action obtained by acting on the first factor, and let
\[
 c_m(t)=\sum_{u\in(\Z/m)^\times}\exp(-2\pi iut/m)
\]
be the integral Ramanujan sum. It satisfies $|c_m(t)|\le\varphi(m)$. Define the integral signed cycle
\[
 Y=\sum_{t=0}^{m-1}c_m(t)\delta(t)_*\mathcal R_0.
\]
Its class is $m$ times the rational primitive-character projector applied to $\cl(\mathcal R_0)$. It lies in the determinant subspace and has all primitive components nonzero, by the preceding nonvanishing for $Q_0$ and injectivity of rational pullback under a finite quotient. The automorphisms preserve the numerical point polarization $H$. Therefore
\begin{align*}
 \operatorname{mass}_H(Y)
 &\le m\varphi(m)\frac{|G_0|^2}{m^{h/2}}\\
 &\le m^2\frac{m^{2h-2}(h!)^2}{m^{h/2}}
   =m^{3h/2}(h!)^2
   \le G^{2g}(g!)^2.
\end{align*}
The extra factor $|G_0|$ makes the integral transfer estimate deliberately generous. No degree estimate for a resolution or an exceptional divisor enters: the relevant cycle is pushed to $W_0$, and its degree is measured against a divisor pulled back from the symmetric product.
\end{proof}

\subsection{A middle-cohomology supplier transfer}

\begin{proposition}[SQ5: keeping the suppliers on curve powers]\label{SQ5}
Choose an identity as in \findingref{SQ3}, and form
\[
 B=\prod_i(C_i'{}^{h_i})^{|n_i|},\qquad
 \ell=\dim B=\sum_i|n_i|h_i\le gR\le L.
\]
There is an integral signed cycle $Y_B$ of dimension $\ell/2$ with
\[
 \operatorname{mass}_{H_B}(Y_B)\le S_B,
\]
where $H_B$ is the sum of the point classes on all its curve factors. For either Weil eigenline $\lambda\subset H^8(A,\C)$, there is a $U$-invariant tensor
\[
 \Gamma\in H^{\ell+8}(B\times A,\C)
\]
whose correspondence action sends $\cl(Y_B)$ to a nonzero vector of $\lambda$. Here $U$ is the centralizer of $\mathcal G$ in the symplectic group of $H^1(C,\Q)$.
\end{proposition}
\begin{proof}
Take the exterior product of the cycles supplied by \findingref{SQ4}, one for every occurrence counted by $|n_i|$. Every $h_i$ is even, so $\ell$ is even. The product-polarization degree formula gives
\begin{align*}
 \operatorname{mass}_{H_B}(Y_B)
 &\le\frac{(\ell/2)!}{\prod_i((h_i/2)!)^{|n_i|}}
       \prod_iS_0^{|n_i|}\\
 &\le(L/2)!S_0^R=S_B.
\end{align*}

The selected determinant line in $H^{h_i}(C_i'{}^{h_i})$ has character $\operatorname{sign}(n_i)a(\chi_i)$. Its Poincar\'e-dual line has the opposite character and also lies in the tensor piece $H^1(C_i')^{\otimes h_i}$. The pairing is nondegenerate: the primitive eigenspaces for $\chi_i$ and $\chi_i^{-1}$ pair perfectly on the curve, and their determinant pairing is perfect. Equivalently, the alternating tensors pair by a nonzero factorial multiple of the determinant of that pairing. Choose $z_i$ in the dual line pairing nontrivially with the corresponding component of $Y_i$.

If $x$ spans the Weil eigenline of character $e_\sigma$, set
\[
 \Gamma=\left(\mathop{\boxtimes}_i
                   z_i^{\boxtimes|n_i|}\right)\boxtimes x.
\]
Its character is zero by the supplier identity. On the order-two character coordinates the absolute values cause no problem, since $|n_i|\equiv n_i\pmod2$. The K\"unneth and projection formulas show that $\Gamma_*\cl(Y_B)$ is $x$ times a product of nonzero pairings. Its cohomological degree is $\ell+8$, so its correspondence action takes $H^\ell(B)$ to $H^8(A)$. Choosing the conjugate lines proves the statement for the other Weil eigenline.

This modifies the supplier transfer of \citet[Proposition \texttt{prop:supply}]{Severn26} by retaining the cycles on curve powers. Their determinant classes are already in middle cohomology. The inverse hard-Lefschetz correspondence used in \citet[Corollary~3.1, p.~25]{Schoen88} is consequently unnecessary for the quantitative argument.
\end{proof}

\subsection{Explicit divisor generators and the mixed-degree estimate}

\begin{proposition}[SQ6: effective output cycles of degree $B_0$]\label{SQ6}
For every cover in the selected Hurwitz component, the eightfold factor $A$ admits finitely many effective integral four-cycles of $\theta_A$-degree at most $B_0$ whose Weil projections span $W_K(A)$. In particular one such cycle has nonzero Weil projection.
\end{proposition}
\begin{proof}
For a quotient $\pi:C\to C'$ of degree $a=|\ker\chi|$, let
\[
 \eta:C'\longrightarrow J(C')\xrightarrow{\pi^*}J(C)
\]
be an Abel--Jacobi map followed by pullback. On rational first cohomology, $\eta^*$ is surjective onto $H^1(C')$ and has a $U$-equivariant right inverse. This follows from the averaging projector for $\ker\chi$ and the identity $\pi_*\pi^*=[a]$. The inclusion $A\hookrightarrow J(C)$ has a rational $U$-equivariant retraction given by $e$. Accordingly the tensor $\Gamma$ of \findingref{SQ5} lifts to a $U$-invariant class on $J(C)^{\ell+1}$ in degree $\ell+8$. Each curve factor occurs in degree one, and the last factor occurs in degree eight, so the first-cohomology splittings give this lift directly in the exterior algebra.

We identify divisors sufficient to span every such invariant. Let $V=H^1(J(C),\Q)$, let $E$ be its polarization form, and let $D$ be the image of $\Q[\mathcal G]$ in $\End(V)$. The algebra $\End_D(V)$ is semisimple and stable under the adjoint involution. Its unitary group is $U$. By \citet[Lemma~3.5]{Milne99}, after scalar extension its unitary elements generate that algebra. Double centralizers then give
\[
 \End_U(V)=D.
\]
The decomposition of $U$ into symplectic, full orthogonal, and paired general linear factors permits application of the classical first fundamental theorems: the invariants in the exterior algebra of any number of copies of $V$ are generated by the invariants in degree two. This is precisely the linear invariant-theory assertion of \citet[Proposition~3.6 and Remark~3.7]{Milne99}. It applies to the algebra $D$ used here and does not require that $D$ be the full endomorphism algebra of the Jacobian.

An invariant alternating form on $V^{\ell+1}$ is $E$ composed with a self-adjoint matrix in $M_{\ell+1}(D)$. Such matrices are spanned by diagonal identities, diagonal $a+a^{-1}$ blocks, and off-diagonal $a,a^{-1}$ pairs, with $a\in\mathcal G$. Their divisor classes are spanned over $\Q$ by
\begin{equation}\label{eq:SQ-divisors}
 \operatorname{pr}_i^*(2\Theta),\qquad
 (\operatorname{pr}_i+a\operatorname{pr}_j)^*(2\Theta),
 \quad a\in\mathcal G,\quad i,j\in\{1,\ldots,\ell+1\},
\end{equation}
including $i=j$. For $i\ne j$, subtracting the two diagonal polarization classes gives the off-diagonal pair. For $i=j$, the identity
$(1+a)^\dagger(1+a)=2+a+a^{-1}$ gives the diagonal block. The harmless factors two do not change the rational span. All the line bundles in \eqref{eq:SQ-divisors} are globally generated because $2\Theta$ is globally generated on an abelian variety. Thus the lift of $\Gamma$ is a complex combination of monomials of
\[
 c=\frac\ell2+4
\]
of these actual divisor classes. No bound on the coefficients of that combination is needed: it will suffice to bound each effective cycle in the finite spanning collection. Notice that the invariant group is $U$, not $\mathcal G$; averaging over $\mathcal G$ would project onto a different space.

We now bound the restrictions of the divisors. Canonical polarizations and the norm identity give
\[
 (\pi^*)^*\theta_{J(C)}=a\theta_{J(C')}
\]
numerically. The Abel--Jacobi curve has degree $g(C')$ for its canonical theta divisor. Riemann--Hurwitz gives
\[
 a(g(C')-1)\le g(C)-1=4G,
\]
and hence
\begin{equation}\label{eq:SQ-curvedegree}
 \deg\eta^*\theta_{J(C)}=a g(C')\le4G+a\le5G.
\end{equation}
On $T=B\times A$, put $H_0=H_B+\theta_A$. The restriction of the product polarization of $J(C)^{\ell+1}$ is at most $5GH_0$ in the nef order. For the divisors in \eqref{eq:SQ-divisors}, the theorem of the square gives a bound of eight times that product polarization. The diagonal case is included because
\[
 (1+a)^*\Theta+(1-a)^*\Theta=4\Theta
\]
numerically. We therefore have the uniform nef inequality
\begin{equation}\label{eq:SQ-nef}
 D_j|_T\le40G H_0.
\end{equation}

For each monomial, choose successive divisors in the globally generated systems so that they contain no component of the preceding effective intersection. This is possible for the finite collection of effective components of the signed cycle $Y_B\times A$. The intersections are effective with integral multiplicities; an empty divisor or a projection of smaller dimension contributes zero. Push the resulting four-dimensional intersections to $A$.

Write $|Y_B|$ for the effective cycle obtained by replacing every coefficient of $Y_B$ by its absolute value. For one monomial the total degree mass of these pushforwards is, using \eqref{eq:SQ-nef}, at most
\begin{align*}
 &(40G)^c\int_{|Y_B|\times A}
                     (H_B+\theta_A)^c\theta_A^4\\
 &\hspace{6mm}=(40G)^c\binom c4
            \operatorname{mass}_{H_B}(Y_B)\int_A\theta_A^8\\
 &\hspace{6mm}\le(40G)^k\binom k4S_B\,8!G^8=B_0.
\end{align*}
Only the term $H_B^{\ell/2}\theta_A^4$ in the expansion survives; this explains the binomial coefficient. The estimate depends on the eightfold dimension, rather than the huge dimension of its ambient Jacobian.

The correspondence of \findingref{SQ5} sends $Y_B$ to a nonzero vector in either Weil eigenline. Since that correspondence is in the complex span of the monomials just considered, the projected classes of the resulting finite collection of effective cycles span both eigenlines. Rationality of the collection and of the Weil projector then gives the assertion over $\Q$. In particular at least one effective cycle in the collection has nonzero projection. Every member has degree at most the total mass bound $B_0$.
\end{proof}

\subsection{An explicit integral-lattice comparison}

\begin{proposition}[SQ7: controlled isogeny to the fixed family]\label{SQ7}
Let $\Lambda$ be an $\cO$-stable lattice of rank eight over $\cO$, with a compatible integral polarization of split Weil type and elementary-divisor product $\Delta$. There is a $K$-isometry between its rational Hermitian space and the standard hyperbolic space such that both the isometry and its inverse have denominators dividing $2\Delta$ relative to suitable integral bases. Consequently the source eightfolds of \findingref{SQ2} can be compared with the fixed target family through isogenies $\phi$ satisfying
\[
 \phi^*\theta_{\rm target}=N_{\rm isog}^{2}\theta_A,
 \qquad N_{\rm isog}=4G^8\Delta.
\]
For effective four-cycles, their pushforward degrees are multiplied by $N_{\rm isog}^{8}$.
\end{proposition}
\begin{proof}
For the alternating form $E$ define
\begin{equation}\label{eq:SQ-integralHermitian}
 h(x,y)=E(x,\omega y)-\bar\omega E(x,y).
\end{equation}
Compatibility with the $\cO$-action shows that $h$ is conjugate-linear in its first argument, linear in its second, and Hermitian. Its values on $\Lambda$ lie in $\cO$. The coefficient of $\omega$ in $h(x,y)$ is $E(x,y)$.

If $H$ is its matrix in an $\cO$-basis, then $|\det H|$ is exactly the product of the elementary divisors of $E$. To check the normalization, the identity Hermitian matrix gives a unimodular alternating form on the integral basis $(v_i,\omega v_i)$. Replacing the identity by $H$ multiplies that alternating determinant by the rational determinant of the $K$-linear matrix $H$, namely
\[
 \Nm_{K/\Q}(\det H)=(\det H)^2.
\]
Thus $\det E=(\det H)^2=\Delta^2$. In signature $(4,4)$, $\det H=\Delta>0$.

Choose a rational isotropic $K$-subspace $I$ of dimension four. The submodule $I\cap\Lambda$ is primitive: if a nonzero $a\in\cO$ satisfies $ax\in I$, then $x\in I$. Since $\cO=\Z[\omega]$ is a principal ideal domain, a basis $v$ of this submodule extends to a basis $(v,w)$ of $\Lambda$. In that basis
\[
 H=\begin{pmatrix}0&B\\B^*&D\end{pmatrix},
 \qquad \Delta=\det H=\Nm(\det B).
\]
Put $b=\det B$ and $X=(B^*)^{-1}D/2$. The substitution $w'=w-vX$ kills the lower-right Hermitian block. Set $u=w'B^{-1}$. Then $h(v,u)=I_4$ and $h(u,u)=h(v,v)=0$, so $(v,u)$ is a standard hyperbolic basis.

Let $\mu$ carry this basis to the standard basis $(e,f)$. The formulas in both directions are
\begin{align*}
 \mu(v)&=e,& \mu(w)&=fB+eX,\\
 \mu^{-1}(e)&=v,& \mu^{-1}(f)&=wB^{-1}-vXB^{-1}.
\end{align*}
The entries of $X$ have denominator $2\bar b$, those of $B^{-1}$ have denominator $b$, and those of $XB^{-1}$ have denominator $2b\bar b=2\Delta$. Multiplication by $2\Delta$ makes both displayed coordinate matrices integral over $\cO$. This proves the two denominator assertions without any bound on the sizes of the entries of $B$ or $D$.

Apply the construction to the source lattice, whose elementary-divisor product $\Delta_A$ divides $G^8$ by \findingref{SQ2}, and to the fixed target lattice, whose product is $\Delta$. Their comparison has denominator dividing $4\Delta_A\Delta$. Hence the integer $N_{\rm isog}=4G^8\Delta$ clears it uniformly.

Transport the rational comparison over a simply connected polydisc in the Hurwitz component. Multiplication by $N_{\rm isog}$ gives a $K$-linear isogeny from each source fibre to the target fibre at the corresponding period point. The unscaled map is a Hermitian isometry, and \eqref{eq:SQ-integralHermitian} recovers the alternating form from it. Therefore
\[
 \phi^*\theta_{\rm target}=N_{\rm isog}^{2}\theta_A.
\]
For a four-cycle $Z$, the projection formula gives
\[
 \deg_{\theta_{\rm target}}\phi_*Z
 =\int_Z\phi^*\theta_{\rm target}^4
 =N_{\rm isog}^{8}\deg_{\theta_A}Z.
\]
Isogeny pushforward preserves nonvanishing of the projected Weil class because it induces a rational isomorphism of Weil planes up to a nonzero scalar.

Finally, the independently checked dominance argument of \citet[Section~9, proof of Theorem~A]{Severn26} gives a nonempty analytic open image in the fixed period space. Applying \findingref{SQ6} and the degree formula on that image proves \findingref{SQ1}.
\end{proof}

\begin{remark}[Integral-action scope]
The source integral action is certified in \findingref{SQ2}. The target hypothesis is satisfied by the fixed principal balanced $E^8$ family and by the integral block-action family of polarization type $(1^4,3^4)$. A raw principal product $E_f^8$ equipped with diagonal scalar $\omega$ and $\bar\omega$ actions, for a nonmaximal CM order of conductor $f>1$, need not have an $\cO_K$-stable lattice. That realization requires an additional conductor denominator and is not silently included in \eqref{eq:SQ-D}. The fixed integral PEL family and its tensor-compatible Hecke orbit satisfy the stated hypothesis.
\end{remark}

\subsection{From an effective seed to numerical bounded division}

\begin{proposition}[SQ8: the integral Weil projector and its degree cost]\label{SQ8}
Let $\delta_4$ and $B_4$ be the eightfold denominator and weighted coefficient sum in Table~\ref{tab:projectors}. Explicitly,
\begin{align*}
 \delta_4&=411891354565851759965785179617340330000,\\
 B_4&=497657009489489182584908205353846834744.
\end{align*}
For an effective four-cycle $Z$ of degree at most $D$, the class $\delta_4P_4\cl(Z)$ has an integral signed representative of degree mass at most $B_4D$. A collection with projected classes spanning the Weil plane therefore gives a spanning collection of integral Weil classes with this mass bound.
\end{proposition}
\begin{proof}
We recall the normalization to make the degree cost explicit. Put $a=\iota(3+\omega)$. Since $a^\dagger a=[7]$, one has $a^*\theta=7\theta$. Its action on $H^8$ has nine distinct eigenvalues
\[
 (3+\omega)^{8-j}(3+\bar\omega)^j,\qquad 0\le j\le8.
\]
The extreme two eigenspaces are exactly the Weil plane. The polynomial of \findingref{P1} is one on those two eigenvalues and zero on the other seven; its cleared numerator is an integral polynomial $Q=\delta_4P_4$. All nine values, the denominator, and all coefficients were recomputed independently in \texttt{severn\_quantitative.py}, and agree with Table~\ref{tab:projectors}.

Pullback of an effective four-cycle by $a^j$ is effective and multiplies its degree by $7^{4j}$. Indeed, $\deg a=7^8$, and the projection formula gives $a_*\theta^4=7^4\theta^4$. Thus the sum of the absolute degree contributions from $Q(a^*)Z$ is at most
\[
 \left(\sum_j|Q_j|7^{4j}\right)D=B_4D.
\]
The class lies in the Weil plane, and multiplication by the nonzero scalar $\delta_4$ does not change the rational span.
\end{proof}

\begin{theorem}[SQ9: an explicit uniform complexity cap]\label{SQ9}
Let $t\ne0$ be a fixed rational flat Weil section on the family of \findingref{SQ1}. Choose $q_t\in\Z_{>0}$ with $q_tt$ integral and put
\[
 T=\int(q_tt)^2\in\Z_{>0},\qquad
 B=B_4D_8(\Delta).
\]
Then on every fibre of the family, and hence on its tensor-compatible Hecke orbit,
\begin{equation}\label{eq:SQ-cap}
 c_t(s)\le C_8(t,\Delta):=B^4(q_t+4T+4).
\end{equation}
In particular there are integral effective cycles with $Nt_s=\cl(Y^+)-\cl(Y^-)$ and both $N$ and the two degrees at most the explicitly specified integer in \eqref{eq:SQ-cap}. The value $t=0$ is represented by the zero cycle with $N=1$.
\end{theorem}
\begin{proof}
On the analytic open subset in \findingref{SQ1}, apply \findingref{SQ8} to its finite spanning collection and select two independent integral Weil classes $y_1,y_2$. Each has a signed representative of mass at most $B$.

Every Weil class is primitive. The polarization has one factor from each $K$-eigenspace, whereas a Weil determinant already fills all eight dimensions of one eigenspace; hence $\theta\wedge y_i=0$. The Hodge star on a primitive $(4,4)$-class in complex dimension eight is $+1$, so the intersection form is positive definite on the Weil plane.

For clarity, the degree-to-norm estimate is given directly on the Weil plane. Choose the invariant K\"ahler metric with class $\theta$ and volume
\[
 V=\int\theta^8/8!=\Delta\ge1.
\]
Let $w$ be a real Weil class of $L^2$-norm one, represented by its invariant harmonic form. Its pointwise norm is $V^{-1/2}$, so its comass is at most $V^{-1/2}$. Wirtinger's identity gives volume $\deg_\theta Z/4!$ for an effective four-cycle $Z$. If $Z_i$ is the signed representative of $y_i$, then $*w=w$ and hence
\[
 |(y_i,w)|=\left|\int_{Z_i}w\right|
 \le\frac{\operatorname{mass}_\theta(Z_i)}{4!\sqrt V}
 \le\frac{B}{4!\sqrt V}.
\]
Taking the supremum over unit vectors $w$ of the positive Weil plane gives $\|y_i\|\le B/(4!\sqrt V)\le B$. This is the direct mass/comass estimate used in \findingref{P0} and \findingref{P3}; no norm assertion about classes outside the Weil plane is needed.

Let $A=((y_i,y_j))$ be the integral Gram matrix. Its entries have absolute value at most $B^2$, and positive definiteness gives an integer
\[
 1\le d:=\det A\le B^4.
\]
The numbers $b_i=(q_tt,y_i)$ are integers and satisfy $|b_i|\le B\sqrt T$. Cramer's rule yields integers $n_1,n_2$ with
\[
 d q_tt=n_1y_1+n_2y_2,
 \qquad |n_1|+|n_2|\le4B^3\sqrt T.
\]
The left multiplier is at most $q_tB^4$, and the signed representative on the right has degree mass at most $4B^4\sqrt T$. Both are bounded by the integer $B^4(q_t+4T+4)$.

This bound holds on a nonempty analytic open subset. The closed-locus assertion of \findingref{R21} implies that $S_{C_8}$ is closed algebraic. Since it contains that analytic open and the base is irreducible, it is the whole base. No Baire selection of an unknown degree is used at this final step: the same explicit integer $C_8$ globalizes.
\end{proof}

\paragraph{A smaller principal-lattice envelope.}
For the principal balanced family, \findingref{P2b} makes the final two-dimensional calculation more concrete. Write $D=D_8(1)$ and set
\begin{equation}\label{eq:SQ-H}
 H=\max\left(1,\left\lfloor
              \frac{\delta_4^2D^2}{31104}\right\rfloor\right).
\end{equation}
For either basis target $u_4$ or $v_4$, one may take the uniform complexity cap
\begin{equation}\label{eq:SQ-principalcap}
 N_s\le1265H,\qquad
 \deg Y_s^\pm\le6325H B_4D.
\end{equation}
For $t=(a u_4+b v_4)/q$, with $a,b\in\Z$ and $q>0$, the corresponding bounds are
\begin{equation}\label{eq:SQ-generaltargetcap}
 N_s\le1265qH,\qquad
 \deg Y_s^\pm\le6325H B_4D(|a|+|b|).
\end{equation}
If a \emph{single common multiplier} is desired, put
\begin{equation}\label{eq:SQ-commonN}
 N_0=1265\operatorname{lcm}(1,\ldots,H).
\end{equation}
Then the multiplier $qN_0$ works on every fibre, with both effective degrees at most
\begin{equation}\label{eq:SQ-commonD}
 5N_0B_4D(|a|+|b|).
\end{equation}

\begin{proof}
Take an effective seed $Z$ on the analytic open and put $z=\delta_4P_4\cl(Z)=xu_4+yv_4\ne0$. By \findingref{P2b},
\[
 \|z\|^2=54(x^2+xy+y^2).
\]
The sharper principal-volume estimate in \findingref{P3} gives
\[
 1\le x^2+xy+y^2\le
       \left\lfloor\frac{\delta_4^2D^2}{54(4!)^2}\right\rfloor\le H.
\]
This explains $31104=54(4!)^2$ in \eqref{eq:SQ-H}.

On the basis $(u_4,v_4)$, $T=\iota(3+\omega)^*$ has matrix
\[
 \begin{pmatrix}-1504&1265\\-1265&-2769\end{pmatrix}.
\]
Consequently
\[
 |\det(z,Tz)|=1265(x^2+xy+y^2).
\]
It is nonzero, and multiplication by it clears the coefficients of either $u_4$ or $v_4$ in the rational basis $(z,Tz)$. Also $\|Tz\|=2401\|z\|$, and the angle $\alpha$ between the two vectors satisfies
\[
 |\sin\alpha|=\frac{1265\sqrt3}{2\cdot2401}.
\]
The shortest nonzero vector of the lattice has norm $\sqrt{54}$, equal to the norms of $u_4$ and $v_4$. Thus, if $u_4=rz+sTz$, or if $v_4=rz+sTz$, elementary two-dimensional Cramer's estimates give
\[
 |r|+2401|s|
 \le\frac{4\cdot2401}{1265\sqrt3}<5.
\]
The signed representatives of $z$ have each degree at most $B_4D$ by \findingref{SQ8}; those of $Tz$ have each degree at most $2401B_4D$. Taking the multiplier $1265(x^2+xy+y^2)$ proves \eqref{eq:SQ-principalcap}. Taking the indicated integral linear combination and multiplying by $q$ proves \eqref{eq:SQ-generaltargetcap}.

Every possible $x^2+xy+y^2$ divides $\operatorname{lcm}(1,\ldots,H)$. Hence $N_0$ in \eqref{eq:SQ-commonN} clears all the basis coefficients simultaneously, proving the fixed-multiplier estimate \eqref{eq:SQ-commonD} on the analytic open. For the complexity caps use the closed $S_C$ of \findingref{R21}; for the common multiplier use its closed fixed-$N$ locus with the degree in \eqref{eq:SQ-commonD}. In both cases the analytic open forces the locus to be the entire irreducible family. The integer $H$ is an explicit upper envelope, not the least attainable quadratic value or denominator.
\end{proof}

\paragraph{Verification scope of the numerical extraction.}
The group order, scalar-unit membership, parabolic basis, cup-product matrix, and all seventeen geometric logarithms were reconstructed exactly before this extraction. The integer-lattice coefficient bound, the integral Schoen transfer, the curve-power contraction, the divisor spanning argument, the $40G$ nef estimate, the binomial intersection formula, the two-sided $2\Delta$ lattice comparison, and the final mass/comass and Gram estimates each received full proof readings, including independent readings by another agent. The projector arithmetic agrees with \findingref{P1}. These are mathematical proof checks and exact arithmetic certificates, not formal theorem-prover verification. The deep cited inputs retain their stated provenance, especially the version-specific supplier repair and dominance theorem in the unrefereed manuscript \citet{Severn26}.
